\documentclass[sigconf]{acmart}

\AtBeginDocument{%
  \providecommand\BibTeX{{Bib\TeX}}}

\setcopyright{none}
\copyrightyear{2026}
\acmYear{2026}
\acmDOI{TO-BE-ASSIGNED}
\acmConference[JCDL '26]{Proceedings of the 2026 ACM/IEEE Joint Conference on Digital Libraries}{October 13--16, 2026}{Dallas, TX, USA}
\acmISBN{TO-BE-ASSIGNED}

\begin{document}

\title[Thevaram 1--8 Corpus Resource]{An Audit-First Public Corpus Resource for Thevaram 1--8}

\author{Bhaveenthan Kajanikanth}
\affiliation{%
  \institution{University of Peradeniya}
  \city{Peradeniya}
  \country{Sri Lanka}}
\email{e22051@eng.pdn.ac.lk}

\author{Uthayasanker Thayasivam}
\affiliation{%
  \institution{University of Moratuwa}
  \city{Moratuwa}
  \country{Sri Lanka}}
\email{rtuthaya@cse.mrt.ac.lk}

\renewcommand{\shortauthors}{Kajanikanth and Thayasivam}

\begin{abstract}
Tamil literary works are increasingly available through digital libraries, but web access alone does not make them reusable research corpora. Computational use needs stable identifiers, source provenance, preserved literary hierarchy, Unicode-safe normalization, and clear separation between poem text and commentary. We present a public GitHub resource for Thevaram books 1--8 derived from Tamil Virtual Academy/TamilVU source pages. The resource contains 8 Thirumurai books, 841 paadal thogupugal, 9012 paadalgal, 9012 commentary records, and 64073 text spans. It includes normalized JSONL tables, schema files, checksums, corpus documentation, rights/access notes, setup commands, and experimental evidence layers for entity annotation and paadal-pozhppurai linking. The Irandaam Thirumurai v1 package is provided as the strongest audited subset, with 122 hymns, 1331 verse records, schema validation, checksums, and audit reports. The contribution is an audit-first resource model for Tamil digital-library research: it makes structure, provenance, quality, and uncertainty visible before downstream retrieval or question answering.
\end{abstract}

\begin{CCSXML}
<ccs2012>
 <concept>
  <concept_id>10002951.10003260.10003282</concept_id>
  <concept_desc>Information systems~Digital libraries and archives</concept_desc>
  <concept_significance>500</concept_significance>
 </concept>
 <concept>
  <concept_id>10010147.10010178.10010179</concept_id>
  <concept_desc>Computing methodologies~Natural language processing</concept_desc>
  <concept_significance>300</concept_significance>
 </concept>
 <concept>
  <concept_id>10010405.10010455.10010458</concept_id>
  <concept_desc>Applied computing~Arts and humanities</concept_desc>
  <concept_significance>300</concept_significance>
 </concept>
</ccs2012>
\end{CCSXML}

\ccsdesc[500]{Information systems~Digital libraries and archives}
\ccsdesc[300]{Computing methodologies~Natural language processing}
\ccsdesc[300]{Applied computing~Arts and humanities}

\keywords{Tamil literature, digital libraries, corpus, Thevaram, provenance, FAIR data, retrieval}

\maketitle

\section{Introduction}

Tamil Virtual Academy/TamilVU provides important public access to Tamil literary and educational material \cite{tamilvu,tamilvulibrary}. Yet a page that can be read in a browser is not automatically a reusable digital-library resource. A retrieval system, corpus analyst, or annotation workflow needs more than page text. It needs source URLs, stable record identifiers, poem boundaries, commentary fields, book and hymn metadata, known gaps, and validation evidence.

This paper presents an audit-first public corpus resource for Thevaram books 1--8, a major Tamil Saiva devotional collection. The resource is designed for Tamil literary research, digital humanities, citation-aware search, and future retrieval-augmented question answering. Its central principle is simple: before building models or interfaces, the corpus must expose provenance, structure, and uncertainty.

Our contributions are:

\begin{enumerate}
  \item a public GitHub resource containing normalized Thevaram 1--8 JSONL tables with 8 books, 841 paadal thogupugal, 9012 paadalgal, 9012 commentary records, and 64073 text spans;
  \item reviewer-facing documentation including schema, checksums, corpus card, datasheet, reuse commands, and rights notes;
  \item an audited Irandaam Thirumurai v1 subset with 122 hymns and 1331 verse records, schema validation, checksums, and audit reports;
  \item experimental evidence layers for entity annotation and commentary linking, explicitly marked as review aids rather than gold expert annotation.
\end{enumerate}

The public resource is available as the \href{https://github.com/bhaveenthank/Bhaveenthan-Tamil-RAG-KnowledgeHub}{Tamil Literary KnowledgeHub GitHub repository}. Project software, schemas, scripts, and original documentation are released under the MIT License. TamilVU-derived source text is attributed to Tamil Virtual Academy/TamilVU and may be subject to separate rights.

\section{Related Work and Resource Gap}

The JCDL Resources Track emphasizes public resources that are reusable, documented, and accessible without review barriers \cite{jcdl2026resources}. Our resource follows the FAIR view that research data should be findable, accessible, interoperable, and reusable \cite{wilkinson2016fair}. We also follow the spirit of datasheets for datasets by documenting composition, collection process, intended uses, and limitations \cite{gebru2021datasheets}.

Language-resource repositories such as LDC-IL show the importance of curated resources for Indian language technology \cite{choudhary2021ldcil}. Cultural-heritage work has also argued that collections should support computational reuse, not only page-level display \cite{ziegler2020opendata}. Tamil scholarly infrastructure is growing as well; for example, Tamilex connects Classical Tamil texts, commentaries, concordances, and lexicographic resources for philological work \cite{tamilex2026}.

The gap addressed here is practical. Tamil literary material exists online, but reusable computational research needs a disciplined data layer. This resource contributes that layer for Thevaram 1--8: normalized tables, source provenance, text spans, validation, release documentation, and evidence layers for future retrieval and annotation.

\section{Resource Design}

Figure~\ref{fig:pipeline} shows the audit-first pipeline. The design separates source inspection, extraction, normalization, validation, release documentation, and downstream evidence layers. This separation is important because literary corpora are easy to flatten. A single page may contain a hymn, poem lines, commentary, source parameters, place names, musical metadata, and navigation context.

\begin{figure}[t]
  \centering
  \fbox{%
    \begin{minipage}{0.94\columnwidth}
      \small
      \textbf{TamilVU pages}
      $\rightarrow$ inspected source URLs
      $\rightarrow$ extraction
      $\rightarrow$ normalized JSONL tables
      $\rightarrow$ schema and checksums
      $\rightarrow$ corpus card and datasheet
      $\rightarrow$ retrieval and annotation evidence
    \end{minipage}}
  \caption{Audit-first resource pipeline. Each stage preserves source provenance and produces reviewable artifacts.}
  \Description{A pipeline from TamilVU pages to inspected source URLs, extraction, normalized JSONL tables, schema and checksums, corpus card and datasheet, and retrieval or annotation evidence.}
  \label{fig:pipeline}
\end{figure}

The core schema uses five families: Thirumurai books, paadal thogupugal, paadal records, commentary records, and text spans. This structure keeps poem text separate from commentary and keeps span-level evidence linked to parent records. The normalized tables are stored as JSONL to make them easy to inspect, stream, validate, and reuse in Python or other tools.

Normalization applies NFC Unicode normalization, removes unsafe invisible/control characters, collapses repeated spaces within lines, trims line edges, and rebuilds span offsets. Poem line boundaries are preserved because line order is part of the literary structure. The repository includes schemas, checksums, and validation reports so that users can check whether a local copy matches the public resource.

\section{Contents and Documentation}

Table~\ref{tab:scope} summarizes the current Thevaram 1--8 normalized resource. These counts describe the full public GitHub resource. The Irandaam Thirumurai v1 package is the most mature audited subset within it.

\begin{table}[t]
  \caption{Thevaram 1--8 normalized corpus scope.}
  \label{tab:scope}
  \centering
  \small
  \begin{tabular}{lrrrr}
    \toprule
    Book & Groups & Poems & Comm. & Spans \\
    \midrule
    1 & 136 & 1469 & 1469 & 10371 \\
    2 & 122 & 1331 & 1331 & 10383 \\
    3 & 126 & 1358 & 1358 & 11837 \\
    4 & 113 & 1070 & 1070 & 6412 \\
    5 & 86 & 1015 & 1015 & 6078 \\
    6 & 99 & 981 & 981 & 9800 \\
    7 & 99 & 1006 & 1006 & 5756 \\
    8 & 60 & 782 & 782 & 3436 \\
    \midrule
    Total & 841 & 9012 & 9012 & 64073 \\
    \bottomrule
  \end{tabular}
\end{table}

The resource includes a top-level README, MIT license for project-owned material, citation metadata, a schema, corpus card, datasheet, usage examples, rights/access statement, and checksum files. A minimal reuse path is:

\begin{verbatim}
python3 -m venv .venv
source .venv/bin/activate
python3 -m pip install -e .
cd data/processed/thevaram_normalized
python3 -m json.tool normalization_summary.json
cd ../../..
shasum -a 256 -c \
  data/releases/thevaram-1-8-v1/corpus_checksum.sha256
\end{verbatim}

The Irandaam Thirumurai v1 subset contains 122 hymns and 1331 verse records. Its audit reports record zero schema violations, zero duplicate song numbers, zero missing verse text records, and complete sampled coverage. This subset is useful for users who want the strictest frozen package before working with the broader 1--8 tables.

\section{Quality Evidence and Reuse}

Table~\ref{tab:review} maps the resource to the JCDL review criteria. This is included to make the resource value explicit.

\begin{table}[t]
  \caption{Resource review criteria and supporting evidence.}
  \label{tab:review}
  \centering
  \small
  \begin{tabular}{p{0.20\columnwidth}p{0.70\columnwidth}}
    \toprule
    Criterion & Evidence \\
    \midrule
    Novelty & Audit-first Tamil literary corpus model with hierarchy, source URLs, poem/commentary separation, checksums, and reviewable evidence layers. \\
    Availability & Public GitHub repository, documentation, setup commands, license, and rights/access statement. \\
    Utility & JSONL tables, schema, corpus card, datasheet, checksums, audit subset, and examples for validation and reuse. \\
    Impact & Supports citation-aware Tamil literary search, corpus analysis, annotation workflows, and future Tamil QA. \\
    \bottomrule
  \end{tabular}
\end{table}

Typical reuse tasks include phrase search over poem text, commentary coverage analysis by book, Tamil lexical or hybrid retrieval, entity review workflows, and citation-grounded question answering experiments. The resource also helps users distinguish evidence types: a poem span, a commentary span, metadata, or a derived annotation are not treated as the same kind of evidence.

Two derived layers are included as experimental evidence. The entity layer contains deterministic mentions linked to registries for deities, temples, nature terms, sacred objects, body parts, and theological concepts. It reports 107761 entity mentions, 624 registry rows, and zero invalid bounds or span mismatches. The paadal-pozhppurai linking layer aligns poem spans to explanatory commentary spans and is best treated as a diagnostic review layer. Many links require review because classical poetic language and prose explanation often do not share surface vocabulary.

\section{Ethics, Rights, AI Use, and Limitations}

The resource is derived from Tamil Virtual Academy/TamilVU pages. It is not an official TamilVU release or mirror. Source text and commentary are attributed to Tamil Virtual Academy/TamilVU. Project-owned code, schemas, scripts, and documentation are MIT licensed; TamilVU-derived source text may be subject to separate rights. Users should check source terms before redistributing source-derived full text.

No human participants or sensitive personal data are involved. AI tools, including OpenAI Codex/ChatGPT, were used to assist with code and script drafting, validation planning, analysis summaries, documentation, and manuscript drafting. The named authors remain responsible for all corpus records, statistics, claims, rights decisions, and final manuscript content.

The resource has limitations. It is not a complete TamilVU corpus and not an expert-verified scholarly edition. Commentary coverage is uneven: book 8 currently has poem-level coverage but no pozhppurai or kurippurai coverage in the normalized commentary table. Entity annotations and paadal-pozhppurai links are review aids, not gold labels. Future work will expand manual review, improve commentary coverage, add retrieval benchmarks, and create a DOI-backed archive when the final rights/access model is settled.

\section{Conclusion}

We presented an audit-first public corpus resource for Thevaram books 1--8. The resource turns TamilVU pages into normalized, source-linked, documented, and checkable corpus tables. Its contribution is not only extracted text, but a reusable method for making Tamil literary data suitable for digital-library research. By preserving provenance, hierarchy, validation evidence, and uncertainty, the resource provides a foundation for citation-grounded Tamil literary search, corpus analysis, annotation review, and future Tamil question answering.

\begin{acks}
The authors thank Tamil Virtual Academy/TamilVU for making Tamil literary materials available online. This paper is not an official TamilVU publication or mirror.
\end{acks}

\bibliographystyle{ACM-Reference-Format}
\bibliography{jcdl2026_thevaram_resource_paper}

\end{document}
